\section{Introduction}

Jiang--Su stability is a regularity property of a \(\mathrm C^*\)-algebra
that also has a dynamical counterpart: an action should absorb the
trivial action on the Jiang--Su algebra up to cocycle conjugacy.
The theorem below establishes this implication for all amenable actions
in the simple, nuclear, finite unital setting. Throughout, tensor products
are minimal, and \(\Zalg\) denotes the Jiang--Su algebra.

\begin{theorem}\label{thm:main}
Let \(A\) be a simple, separable, unital, infinite-dimensional, nuclear,
stably finite complex \(\mathrm C^*\)-algebra satisfying
\(A\cong A\otimes\Zalg\). Let \(G\) be a countable discrete amenable group
and let \(\alpha:G\to\Aut(A)\) be any action. Put
\[
 B=A\otimes\Zalg,\qquad
 \beta_g=\alpha_g\otimes\operatorname{id}_{\Zalg}.
\]
There are a unital \(*\)-isomorphism \(\Phi:A\to B\) and unitaries
\(u_g\in B\), \(g\in G\), such that
\begin{align}
 u_e&=1,& u_{gh}&=u_g\beta_g(u_h),\label{eq:cocycle}\\
 \Phi(\alpha_g(a))&=u_g\beta_g(\Phi(a))u_g^*
 &&(g\in G,\ a\in A).\label{eq:conjugacy}
\end{align}
\end{theorem}

Here stable finiteness means that every isometry in every matrix algebra
over \(A\) is unitary. Amenability is used in its F{\o}lner form: for
every finite \(H\subset G\) and \(\delta>0\), some nonempty finite
\(F\subset G\) satisfies
\(\lvert hF\mathbin{\triangle}F\rvert<\delta\lvert F\rvert\) for \(h\in H\).
Theorem~\ref{thm:main} imposes no outerness, Rokhlin, or UCT assumption,
and no condition on the extreme tracial boundary or on trace orbits.
In particular, it includes finite groups and trivial actions.

\subsection{Context and scope}

Szab\'o formulated automatic equivariant Jiang--Su stability for
amenable actions on separable simple nuclear Jiang--Su-stable algebras
as Conjecture~A of \cite{SzaboSI2021}.
Theorem~\ref{thm:main} resolves its unital, stably finite case.
The original formulation also includes nonunital algebras, which lie
beyond the scope of the present theorem. Earlier affirmative cases
include finite algebras with finitely many rays of extremal traces
and very weak comparison \cite[Theorem~C]{SzaboSI2021}.

Szab\'o and Wouters proved that equivariant uniform property Gamma
characterizes equivariant Jiang--Su stability in the setting used here
\cite[Corollary~C]{SzaboWouters2025}. Their theorem supplies the
absorption step of the proof. Our task is to construct the required
invariant central trace splittings for an arbitrary action.
These splittings are uniform over the entire trace space, even when
the action moves traces.

\subsection{Construction and proof structure}

The technical argument has three parts.
First, tensorial absorption supplies countably many commuting
\emph{orbit algebras} inside the uniform tracial central sequence algebra.
Under every extreme trace on the algebra they generate, distinct orbit
algebras are independent. This intrinsic independence is useful because
the traces under consideration need not extend to the ambient algebra.
We then prove a joint central limit statement for normalized sums and
their empirical covariances, allowing the extreme trace to vary with
the number of rows.

Second, the covariance is incorporated into a quadratic Gaussian
construction. For a F{\o}lner set \(F\), its Gram vectors define a
positive operator \(P_F\) of trace one. Taking a polynomial approximation
to \(P_F^{1/2}\) gives a quadratic form whose variance stays bounded
away from zero and whose change under translation is small.
This square-root normalization handles covariance matrices of arbitrary
rank, including highly correlated ones. Empirical covariances turn the
same formula into exactly equivariant noncommutative polynomials.
Compactness of the finite covariance space then promotes the scalar
moment limits to estimates uniform over every tracial state.

Finally, weighting traces by elements of \(A\) yields a central invariant
unitary with Haar distribution relative to \(A\). Two explicit
diagonal selections give this unitary and then equal trace projections.
The construction of the rows appears in Section~\ref{sec:rows}, the
uniform unitary estimates in Section~\ref{sec:gaussian}, and the
diagonal selections and absorption argument in Section~\ref{sec:splittings}.
All parts of this construction are proved below; the final absorption
theorem is stated separately in Theorem~\ref{thm:absorption}.

Several ingredients can be used separately under their stated hypotheses:
intrinsic independence for commuting rows
(Lemma~\ref{lem:factor-independence}), moment estimates under varying
tracial states (Proposition~\ref{prop:joint-clt}), and convergence of
scalar spectral laws against bounded continuous functions without
uniform operator-norm bounds (Lemma~\ref{lem:quadratic-moment-transfer}).
